% impulse_ft: impulse is the area under the force-time curve (Ch 7, fig:impulse)
% Two force pulses deliver the same impulse (the same change in momentum).
% Each is a Gaussian F = F0 exp(-(t-t0)^2/(2 s^2)), whose area is F0 s sqrt(2 pi);
% equal areas therefore need F0 s to be the same for both.
%   hard stop: s = 0.25, F0 = 1       airbag: s = 1.10, F0 = 0.25/1.10 = 0.227
% The airbag spreads the impulse over 4.4 times as long, so its peak force is
% 4.4 times smaller.  (The real ratio is larger; it is compressed here to fit.)
\begin{tikzpicture}
\begin{axis}[bookplot,
    scale only axis, width=8.3cm, height=4.3cm,
    xmin=0, xmax=8.4, ymin=0, ymax=1.2,
    xlabel={time $t$}, ylabel={force $F$},
    xlabel style={at={(axis description cs:1,-0.02)},anchor=north east},
    ylabel style={at={(axis description cs:-0.02,0.5)}}]
  \def\sa{0.25}\def\ta{1.5}           % hard stop
  \def\sb{1.10}\def\tb{4.7}           % airbag
  % shaded areas (the impulses)
  \addplot[draw=none,fill=cAcc,fill opacity=0.18,domain=0:3,samples=100]
    {exp(-(x-\ta)^2/(2*\sa^2))} \closedcycle;
  \addplot[draw=none,fill=cPrim,fill opacity=0.18,domain=0.2:8.2,samples=120]
    {(\sa/\sb)*exp(-(x-\tb)^2/(2*\sb^2))} \closedcycle;
  % the two pulses
  \addplot[cAcc,domain=0:3,samples=100] {exp(-(x-\ta)^2/(2*\sa^2))};
  \addplot[cPrim,domain=0:8.2,samples=120] {(\sa/\sb)*exp(-(x-\tb)^2/(2*\sb^2))};
  % labels
  \node[cAcc,anchor=west,align=left] at (axis cs:1.9,0.86)
    {hard stop:\\ brief, large force};
  \node[cPrim,anchor=south,align=center] at (axis cs:\tb,0.30)
    {airbag:\\ longer, gentler};
  \node[note,anchor=north east,align=right] at (axis cs:8.4,1.17)
    {equal areas: the same impulse,\\ the same $\Delta p$};
\end{axis}
\end{tikzpicture}
